\documentclass[11pt]{article}

% Change "review" to "final" for the camera-ready version.
\usepackage[review]{acl}

\usepackage{times}
\usepackage{latexsym}
\usepackage[T1]{fontenc}
\usepackage[utf8]{inputenc}
\usepackage{microtype}
\usepackage{inconsolata}
\usepackage{graphicx}
\usepackage{booktabs}
\usepackage{makecell}

\title{Can Prompt Engineering Bridge the Long-Context NLI Reasoning Gap?}

\author{
  Your Name \\
  Your Department \\
  Your University \\
  \texttt{your@email} \\
  \And
  Teammate Name \\
  Your Department \\
  Your University \\
  \texttt{teammate@email} \\
}

\begin{document}
\maketitle

\begin{abstract}
% TODO: 3-5 sentences -- the problem, why it matters, your approach,
% and the headline finding.
\end{abstract}

\section{Introduction}

Natural language inference (NLI) asks whether a hypothesis is entailed by,
contradicted by, or neither supported nor contradicted by a premise, and is
widely used to evaluate language understanding. However, influential
benchmarks such as SNLI and MultiNLI predominantly consist of sentence-level
premise--hypothesis pairs \citep{bowman-etal-2015-large,williams-etal-2018-broad}.
High accuracy on these datasets does not necessarily indicate robust
reasoning, as models can exploit annotation artifacts and shallow syntactic
heuristics that correlate with the labels
\citep{gururangan-etal-2018-annotation,mccoy-etal-2019-right}. Evaluating NLI
over longer passages therefore provides a stricter test, requiring models to
locate and integrate evidence distributed across multiple sentences before
assigning a label.

ConTRoL addresses this limitation by evaluating NLI over long, expert-designed
passages that require reasoning across multiple sentences and paragraphs. The
dataset contains 8,325 premise--hypothesis pairs, with premises averaging 452
words, and covers logical, information-integration, coreferential, temporal,
and analytical reasoning. In the original study, the strongest fine-tuned
baseline achieved 60.95\% accuracy, compared with a human ceiling of 94.4\%,
with particularly low performance on logical reasoning and information
integration. ConTRoL therefore tests more than the ability to process a long
input: models must identify, combine, and interpret the evidence relevant to
the hypothesis \citep{liu-etal-2021-natural}.

Modern instruction-tuned language models enable task adaptation without
ConTRoL-specific fine-tuning through direct instructions, in-context
demonstrations, and intermediate evidence-seeking steps. Prior work has shown
that demonstrations can guide model behavior without parameter updates
\citep{brown2020language}, while question-based methods can support factual
verification and multi-step reasoning by decomposing complex claims and
gathering relevant evidence \citep{wang-etal-2020-asking,min-etal-2019-multi}.
However, the relative contributions of these mechanisms remain unclear for
NLI over long passages with three possible labels. We therefore ask whether
question-guided evidence extraction improves performance beyond direct
prompting and the retrieval of relevant premise sentences alone.

To answer this question, we evaluate three modern open-weight language models
on ConTRoL using six strategies that differ in how they retrieve and
structure evidence. Two single-call baselines use zero-shot and three-shot
prompting, while a retrieve-then-classify control isolates the contribution
of locating relevant premise sentences. Three question-guided pipelines
generate evidence-seeking questions from the hypothesis alone, the premise
alone, or both the premise and hypothesis, before locating supporting
sentences, extracting answers, aggregating the resulting evidence, and
predicting the NLI label. This design separates the contributions of
demonstrations, evidence retrieval, question generation, and the information
available during decomposition, while also enabling analysis of the
pipelines' intermediate questions and extracted evidence. Figure~\ref{fig:example}
shows a worked example of this process on a real ConTRoL instance.

\begin{figure}[t]
\centering
\small
\fbox{%
\begin{minipage}{0.94\linewidth}
\textbf{Premise:} ``Most of the private companies have decided against
awarding annual increase in the salaries of their employees for the previous
year due to the current economic situation.''

\textbf{Hypothesis:} These companies may announce a hike in salaries next
year.

\vspace{4pt}
\textbf{Q1 (bridge/comparison):} Does the premise mention any plans or
possibilities for salary increases in the future?\\
\textbf{A1:} The premise mentions that private companies have decided
against salary increases for the previous year, but does not provide
information about future plans or possibilities for salary increases.

\vspace{4pt}
\textbf{Q2 (regular):} What is the reason given in the premise for not
awarding annual salary increases this year?\\
\textbf{A2:} The reason given is the current economic situation.

\vspace{4pt}
\textbf{Predicted label:} Neutral (matches gold) --- the premise explains
why last year's raise was withheld, but says nothing about next year, so it
neither supports nor contradicts the hypothesis.
\end{minipage}%
}
\caption{A \texttt{bridge-question} example from ConTRoL (\texttt{id\_4089}):
questions are generated jointly from the premise and hypothesis, answered
against the premise, and passed to the classifier along with the hypothesis.
Q1 tests the hypothesis's claim directly; Q2 surfaces supporting context.
Because the premise never addresses next year, the classifier correctly
predicts Neutral rather than guessing.}
\label{fig:example}
\end{figure}

\section{Related Work}

Popular sentence-level NLI benchmarks such as SNLI and MultiNLI
\citep{bowman-etal-2015-large,williams-etal-2018-broad} have a known problem:
models can score well on them by learning shortcuts and surface patterns
instead of actually reasoning about the premise and hypothesis
\citep{gururangan-etal-2018-annotation,mccoy-etal-2019-right}. Most fixes for
this happen during training, for example by filtering out easy examples or
adding extra supervision. ConTRoL takes a different approach: it makes these
shortcuts harder to use by requiring evidence to be found and combined across
a long, multi-sentence premise \citep{liu-etal-2021-natural}. We take yet
another approach. Instead of retraining or re-annotating data, we test
whether prompting and structured question generation, used only at inference
time with no fine-tuning at all, can close this performance gap.

Working only at inference time means we have to rely on prompting to adapt
the model to the task. In zero-shot prompting, the task is described using
instructions alone. Few-shot in-context learning adds example demonstrations
that show the model what the labels mean \citep{brown2020language}.
Chain-of-thought prompting goes a step further, adding worked-out reasoning
steps to these examples \citep{wei2022chain}. Our zero-shot and three-shot
conditions test only the effect of adding demonstrations: since our
three-shot examples do not include any reasoning steps, we treat them as
plain in-context learning, not chain-of-thought.

Beyond demonstrations, our question-based methods build on a different line
of work: using generated questions to check and structure evidence, instead
of just showing the model labeled examples. QA and NLI are closely linked
tasks: \citet{demszky2018transforming} show that QA datasets can be
mechanically turned into NLI examples, by rewriting each question-answer
pair as a declarative hypothesis paired with the source passage as premise.
Our pipelines lean on the reverse direction of that connection, using
generated questions and their answers as intermediate evidence for an NLI
classifier, rather than converting them into NLI training data directly.
QAGS and FEQA generate questions
from a summary and answer them using the source document, to check whether
the summary stays factually consistent with it
\citep{wang-etal-2020-asking,durmus-etal-2020-feqa}. FActScore instead breaks
long generated text into small, independently checkable facts
\citep{min-etal-2023-factscore}. Multi-hop QA work such as HotpotQA and
DecompRC shows the value of a related idea: finding supporting evidence and
breaking a complex question into simpler ones
\citep{yang-etal-2018-hotpotqa,min-etal-2019-multi}. None of these methods
target three-way inference over long passages; they check summaries or
answer single questions instead. We adapt their core idea to passage-level
NLI: we generate questions from the hypothesis, the premise, or both, and
add a retrieval-only control to show how much of the benefit simply comes
from finding the right sentences, versus from the generated questions
themselves.

\section{Method}
\label{sec:method}

\subsection{Dataset}
\label{sec:dataset}

We evaluate on \textbf{ConTRoL} \citep{liu-etal-2021-natural}, a dataset for
document-level natural language inference constructed from expert-designed
reading-comprehension material. Each ConTRoL instance pairs a long premise
(averaging approximately 500 words) with a short hypothesis, and requires
identifying whether the hypothesis is \textbf{Entailed}, \textbf{Contradicted},
or \textbf{Neutral} with respect to the premise. The dataset comprises 8,325
premise-hypothesis pairs. Premises are long, multi-sentence passages rather
than single sentences, so the evidence needed to decide a label is often
spread across several sentences rather than contained in one. This makes
ConTRoL more demanding than sentence-level NLI benchmarks.

\subsection{Sampling}
\label{sec:sampling}

All experiments were run against fixed samples drawn from the ConTRoL train
split, using a fixed random seed (seed~$=$~42).\footnote{A na\"ive
lexicographic sort over string-formatted instance IDs (e.g.\ \texttt{id\_1},
\texttt{id\_10}, \texttt{id\_100}, \ldots, \texttt{id\_11}) does not
correspond to a numeric or random ordering and yields a non-representative
sample.} The sampled instance IDs were persisted once and reused identically
across every model, ensuring that all models and methods are evaluated on the
exact same data.

Due to the substantially higher computational cost of the premise-decomposition
methods (Section~\ref{sec:methods}), we adopt two sample sizes: $N=1{,}500$
for the six methods that issue a bounded number of model calls per instance,
and $N=150$ for the three premise-question configurations, whose per-instance
cost scales with the number of decomposed sub-questions.

\subsection{Seeding Strategies}
\label{sec:seeding}

Two of our pipelines (Section~\ref{sec:methods}) generate probe questions from
an anchor set of extracted spans, rather than directly from the raw text. We
compare two seeding strategies for producing this anchor set.

\paragraph{Part-of-Speech (POS) seeding.} We use NLTK's POS tagger
\citep{loper2002nltk} to extract noun phrases, main verbs, and numeric spans
from the source text via shallow chunking. This method requires no model
inference and is fully deterministic.

\paragraph{Semantic Role Labeling (SRL) seeding.} Following the QA-SRL
paradigm \citep{fitzgerald2018large,pyatkin2021asking}, we prompt the
evaluated LLM to extract PropBank-style predicate-argument structures ---
agent (ARG0), theme (ARG1), recipient (ARG2), and modifier roles for time,
location, manner, and cause --- for each predicate in the source text. Each
role-argument pair is converted into an anchor of the form
\texttt{\{role\}: \{span\} [predicate: \{verb\}]}. This strategy is more
computationally expensive than POS seeding, requiring one additional model
call, but yields anchors with explicit predicate-argument structure rather
than flat lexical spans.

\subsection{Methods}
\label{sec:methods}

We evaluate nine method configurations, organized into three groups.

\paragraph{Direct baselines.} \texttt{zero-shot} issues the premise and
hypothesis to the model in a single call with no intermediate reasoning step.
\texttt{few-shot-CoT} uses the same single-call structure, prepending three
fixed demonstrations --- one worked example per label, held constant across
all evaluations --- following \citet{wei2022chain}.

\paragraph{Retrieval control.} \texttt{retrieve-then-classify} isolates the
contribution of evidence localization from that of question generation. A
locator model identifies the premise sentences most relevant to the
hypothesis, using the \textbf{hypothesis itself} as the retrieval query; the
classifier is then given these sentences verbatim, without an intermediate
answer-extraction step. This condition allows us to attribute any additional
gain from the question-based methods below specifically to the generated
questions, rather than to evidence retrieval alone.

\paragraph{Question-based pipelines.} The remaining six configurations share a
common three-stage architecture: (i) question generation, (ii) evidence
localization and answer extraction, and (iii) classification of the
hypothesis given the resulting (question, answer) pairs. Stages (ii) and
(iii) are implemented identically across all six configurations (detailed
below), so that any difference in downstream accuracy can be attributed to
\emph{which} questions were generated and \emph{which text they were
generated from}, not to how they are answered or judged. The three
question-based methods differ precisely along this second axis --- questions
generated from the hypothesis, from the premise, or from both --- letting us
compare which side of the premise-hypothesis pair is more productive to
interrogate.

\subsubsection{Answering and Classification}

\paragraph{Locate and answer.} A locator model reads the premise and one
question and returns up to five relevant sentence indices (not necessarily
adjacent, allowing multi-hop evidence). An extractor model then answers the
question using only those sentences, returning a not-answerable marker
instead of guessing if they do not address it; such questions are dropped.

\paragraph{Classify.} The surviving (question, answer) pairs and the
hypothesis are given to the classifier in one call, which returns a single
Entailment / Contradiction / Neutral label. Priority is Contradiction over
Entailment over Neutral, and one clearly decisive piece of evidence is
treated as sufficient on its own, without being diluted by other, merely
tangential evidence in the same call.

Stage (i) --- question generation --- is the only stage that differs between
the three methods below.

\begin{itemize}
    \item \textbf{h-question} generates 2--4 probe questions directly from
    the hypothesis, seeded via either POS or SRL extraction
    (Section~\ref{sec:seeding}) applied to the hypothesis text. This pipeline
    is hypothesis-aware but premise-blind at generation time.

    \item \textbf{bridge-question} generates 2--4 sub-questions from the
    premise and hypothesis jointly, with at least one question explicitly
    required to test the hypothesis's core claim against premise evidence,
    following the bridging/comparison framing of multi-hop QA
    \citep{yang-etal-2018-hotpotqa,min-etal-2019-multi}. This is the only configuration
    whose question-generation stage observes both texts simultaneously.

    \item \textbf{p-question (decomposition)} applies \emph{atomic
    decomposition} to the premise: the model is prompted to break the
    premise into the smallest standalone factual claims it contains ---
    each covering exactly one date, name, number, or event --- following
    the atomic-fact paradigm of FActScore \citep{min-etal-2023-factscore}. A claim
    is ``atomic'' in the sense that it cannot be split further without
    losing a checkable unit of information; a premise stating that ``the
    government paid \$1~lakh to families of railway-accident victims''
    atomizes into separate facts for the amount, the recipients, and the
    triggering event. In a second step, the model additionally identifies
    \emph{relations} between these facts --- comparisons, causal links, and
    temporal orderings --- following relation-level verification in the
    style of \citet{goyal2020evaluating,goyal2021annotating}. Each generated
    question is tagged as targeting either a single fact or a relation
    between facts. This two-step decomposition guarantees that both the
    premise's individual claims and the connections between them are
    covered by construction, rather than left to the model's unconstrained
    judgment of what is salient. This mode is premise-aware but
    hypothesis-blind at generation time.

    \item \textbf{p-question (seeded)} generates one question per anchor
    extracted from the premise via the seeding strategies of
    Section~\ref{sec:seeding}, evaluated separately for POS and SRL
    seeding.
\end{itemize}

Stage-1 question generation additionally receives three worked examples
in-context for all six question-based configurations, following the same
few-shot design as \texttt{few-shot-CoT}.

\subsection{Models}
\label{sec:models}

We evaluate four models spanning both self-hosted open-weight and API-hosted
settings: \textbf{Qwen2.5-32B} and \textbf{gpt-oss-20B}, run locally via
Ollama-backed infrastructure, and \textbf{GPT-4o-mini} and
\textbf{DeepSeek-V4-Flash}, accessed via their respective provider APIs. This
selection allows comparison across model scale and serving infrastructure,
while controlling the evaluation pipeline, prompts, and sampled data
identically across all four.


\section{Experimental Evaluation}

\subsection{Experimental Setting}

We evaluate five models spanning self-hosted and API-hosted settings
(Section 3.5): Qwen2.5-32B, gpt-oss-20B, GPT-4o-mini, DeepSeek-V4-Flash, and
Gemma-3-27B. Following Section 3.2, the six non-premise methods are evaluated
on N=1,500 fixed-seed samples from the ConTRoL train split, identical across
all models; the three p-question methods are evaluated on the first N=150 of
that same sample, due to their substantially higher per-instance cost. We
report accuracy as the primary metric, alongside macro-averaged
precision/recall/F1 (robust to label imbalance) and the Matthews Correlation
Coefficient, a single robust summary statistic standard in the GLUE benchmark
\citep{wang2018glue}. At the time of writing, gpt-oss-20B and GPT-4o-mini are
missing bridge-question and, for GPT-4o-mini, h-question (srl) and the
p-question trio, due to provider-side rate limits encountered during the
production run; these cells are marked `--' below.

\subsection{Key Results}

\begin{table}[t]
\centering
\small
\renewcommand{\arraystretch}{0.92}
\setlength{\tabcolsep}{3.5pt}
\begin{tabular}{lcccccc}
\toprule
Model & \makecell{zero-\\shot} & \makecell{few-\\shot} & \makecell{retrieve-\\classify} & \makecell{h-quest.\\(pos)} & \makecell{h-quest.\\(srl)} & \makecell{bridge-\\question} \\
\midrule
Qwen2.5-32B & 67.1 & 69.6 & 69.4 & 67.1 & 66.0 & 69.0 \\
gpt-oss-20B & 70.8 & 70.0 & 68.2 & 64.7 & 65.9 & -- \\
GPT-4o-mini & 65.9 & 65.5 & 67.1 & 59.8 & -- & -- \\
DeepSeek-V4-Flash & 74.6 & 73.5 & 71.3 & 59.4 & 69.3 & 71.9 \\
Gemma-3-27B & 64.5 & 64.9 & 63.5 & 64.5 & 63.2 & 67.0 \\
\bottomrule
\end{tabular}
\caption{Accuracy (\%) across the six N=1,500 methods and five models. `--'
marks gpt-oss-20B's not-yet-completed bridge-question run and GPT-4o-mini's
pending bridge-question/h-question~(srl) reruns (Section 4.1).}
\label{tab:main-results}
\end{table}

Table~\ref{tab:main-results} reports accuracy across the six N=1,500 methods
and five models. Two patterns stand out. First, no method uniformly wins:
DeepSeek-V4-Flash's strongest results come from its baselines (zero-shot
74.6\%, few-shot 73.5\%), which outscore every question-based method on that
model, including bridge-question (71.9\%). This pattern recurs, to varying
degrees, across all five models --- question-guided evidence extraction does
not reliably beat direct prompting or retrieval-only classification in this
evaluation. Second, within the question-based family, bridge-question is
consistently the strongest configuration, consistent with it being the only
method whose question-generation stage observes both the premise and
hypothesis jointly.

\subsection{Premise-Question Methods: An Independent Comparison}

\begin{table*}[t]
\centering
\small
\renewcommand{\arraystretch}{0.92}
\setlength{\tabcolsep}{3.5pt}
\begin{tabular}{lccccc|ccc}
\toprule
Model & \makecell{zero-\\shot} & \makecell{few-\\shot} & \makecell{retrieve-\\classify} & \makecell{h-quest.\\(pos)} & \makecell{bridge-\\question} & \makecell{p-quest.\\(decomp)} & \makecell{p-quest.\\(seed,pos)} & \makecell{p-quest.\\(seed,srl)} \\
\midrule
Qwen2.5-32B & 70.0 & 74.7 & 78.0$^\dagger$ & 74.7 & 73.3 & 59.3 & 61.3$^\dagger$ & 60.7 \\
DeepSeek-V4-Flash & 80.0$^\dagger$ & 80.0 & 79.3 & 71.3 & 77.3 & 68.7$^\dagger$ & 62.7 & 62.0 \\
Gemma-3-27B & 68.0 & 68.7 & 64.0 & 68.0 & 74.0$^\dagger$ & 55.3$^\dagger$ & 47.3 & 51.3 \\
\bottomrule
\end{tabular}
\caption{Accuracy (\%) on the identical 150-item subset used by the
p-question methods, letting every method be compared fairly against
p-question despite its smaller N (Section 3.2). $^\dagger$ marks the
per-model best-baseline/best-p-question pair tested with a McNemar exact
test \citep{dror-etal-2018-hitchhikers}; all three are significant at
$p<0.01$.}
\label{tab:pq-fair}
\end{table*}

The three p-question variants were evaluated at N=150 rather than N=1,500
(Section 3.2), so comparing raw numbers risks conflating sample and method
differences. Table~\ref{tab:pq-fair} instead restricts every method to the
identical 150 items, giving a true paired comparison. The gap widens, not
narrows: each model's best baseline scores higher on this slice than its
full-1,500 average (e.g. DeepSeek-V4-Flash's zero-shot reaches 80.0\% vs.
74.6\% overall), while p-question stays roughly flat. A McNemar exact test
\citep{dror-etal-2018-hitchhikers} on each model's
best-baseline-vs-best-p-question pair confirms this is not noise: all three
differences are significant at $p<0.01$.

The bottleneck is generation, not answering. Restricted to questions Stage 1
actually produced, the Locator/Extractor answers 96-100\% of them across
every model and variant. But 21-79\% of samples never got a usable question
from Stage 1, forcing a zero-shot or hardcoded-Neutral fallback.
p-question's weakness is that premise-only decomposition often generates
nothing to answer, corroborating an earlier manual trace where 3 of 4
checked failures were missing questions, not bad answers.

Within p-question, no seeding strategy dominates uniformly: decomposition is
strongest for two of three models, seeded-pos for the third.

\section{Conclusion}
% TODO

\section*{Limitations}
% TODO -- required section. The N=1500 vs N=150 sample-size asymmetry
% (Section~\ref{sec:sampling}) belongs here at minimum.

\bibliography{references}
\bibliographystyle{acl_natbib}

\end{document}
